\documentclass[11pt]{article}
\usepackage[margin=1in]{geometry}
\usepackage[T1]{fontenc}
\usepackage[utf8]{inputenc}
\usepackage{amsmath,amssymb,amsthm}
\usepackage{graphicx}
\usepackage{xcolor}
\usepackage{microtype}
\usepackage[colorlinks=true,linkcolor=blue!55!black,citecolor=blue!55!black,urlcolor=blue!55!black]{hyperref}

\newtheorem{theorem}{Theorem}
\newtheorem{lemma}[theorem]{Lemma}
\newtheorem{corollary}[theorem]{Corollary}
\newcommand{\T}{\mathbb T}
\newcommand{\Z}{\mathbb Z}
\newcommand{\R}{\mathbb R}
\newcommand{\F}{\mathcal F}

\title{The Universal Tiling Conjecture for periods at most five}
\author{Candidate partial result for arXiv:1209.0666}
\date{12 August 2026}

\begin{document}
\maketitle

\begin{center}
\fbox{\parbox{0.92\textwidth}{\raggedright
\textbf{Status: candidate substantial partial result, likely valid.}
The Universal Tiling Conjecture $\mathrm{UTC}(p)$ of Dutkay and
Jorgensen holds for every $1\le p\le5$.  Hence every bounded measurable
spectral set of measure one with a spectrum of integer period at most five
tiles the line.  The arbitrary-period conjecture remains open.}}
\end{center}

\section{Source problem}

Dutkay and Jorgensen~\cite{source} formulate the following conjecture on
PDF page~2 (Conjecture~1.4).  Let
\(\Gamma=\{\lambda_0=0,\lambda_1,\ldots,\lambda_{p-1}\}\subset\R\)
have $p$ elements and at least one spectrum of the form $p^{-1}A$ with
$A\subset\Z$.  If $A_1,\ldots,A_m\subset\Z$ are any finite family such
that $p^{-1}A_i$ is a spectrum for $\Gamma$, must there be one set
$\mathcal T\subset\Z$ for which
\[
 A_i\oplus\mathcal T=\Z\qquad(1\le i\le m)?
 \tag{UTC($p$)}
\]
Here $\oplus$ means that every integer has a unique representation.

\begin{figure}[ht]
\centering
\includegraphics[width=\textwidth]{figures/open_problem_crop.png}
\caption{Conjecture~1.4 in arXiv:1209.0666, source PDF page~2.}
\end{figure}

\section{Idea of the proof}

Encode the integer frequency $a$ by the unimodular row vector
\[
 v(a)=\left(e^{2\pi i a\lambda_j/p}\right)_{j=0}^{p-1}\in\T^p.
 \tag{1}
\]
Then $p^{-1}A$ is a spectrum for $\Gamma$ exactly when the vectors
$p^{-1/2}v(a)$, $a\in A$, are an orthonormal basis of $\mathbb C^p$.
Moreover, $v(a+b)=v(a)v(b)$ coordinatewise.  Thus the problem concerns
Hadamard bases cut out of one cyclic character orbit.

For $p=2,3$, a zero sum of $p$ unit vectors is rigid: it is an antipodal
pair or an equilateral triple.  A gcd then puts every $A_i$ on the same
complete residue scale.  For $p=4$, every dephased Hadamard matrix has
rows $1,u,w,uw$ with $u^2=1$; the remaining two-point quotient bases have
one fixed 2-adic scale.  For $p=5$, the unique Hadamard equivalence class
is the Fourier matrix, so every basis is the unique order-five subgroup
of the cyclic orbit.  These descriptions exhibit a common complement.

\section{The low-period theorem}

\begin{theorem}\label{main}
The Universal Tiling Conjecture $\mathrm{UTC}(p)$ holds for
\[
 p\in\{1,2,3,4,5\}.
\]
More explicitly, for every admissible $\Gamma$ and every finite family of
integer spectra $p^{-1}A_i$ of $\Gamma$, the proof below constructs one
periodic set $\mathcal T\subset\Z$ such that
$A_i\oplus\mathcal T=\Z$ for every $i$.
\end{theorem}

\begin{proof}
The case $p=1$ is immediate.  We use the encoding (1) throughout.

\medskip
\noindent\textbf{The cases $p=2,3$.}
Let $D$ be the finite set of all nonzero differences $a-b$ with
$a,b\in A_i$ for some $i$, and put
\[
 g=\gcd\{|d|:d\in D\}.
\]
For each $d\in D$, orthogonality gives
\[
 \sum_{j=0}^{p-1}e^{2\pi i d\lambda_j/p}=0.
 \tag{2}
\]

If $p=2$, (2) says
$e^{\pi i d\lambda_1}=-1$, so $d\lambda_1$ is an odd integer.  Since
$g$ is an integral linear combination of the elements of $D$,
$c=g\lambda_1$ is an integer.  It must be odd, and then
$(d/g)c=d\lambda_1$ odd implies that $d/g$ is odd for every $d\in D$.

If $p=3$, the elementary equality
\[
 1+z+w=0,\qquad |z|=|w|=1,
\]
forces $\{z,w\}=\{e^{2\pi i/3},e^{4\pi i/3}\}$.  Therefore
$d\lambda_1$ and $d\lambda_2$ are integers not divisible by $3$.
Again $c=g\lambda_1\in\Z$, now $3\nmid c$, and hence
$3\nmid d/g$ for every $d\in D$.

In either case, the elements of each $A_i$ are congruent modulo $g$, and
after subtracting one element and dividing by $g$, their $p$ residue
classes modulo $p$ are pairwise distinct.  They are therefore all the
classes $0,1,\ldots,p-1$.  It follows directly, or by reducing modulo
$pg$, that the single set
\[
 \mathcal T=\{0,1,\ldots,g-1\}+pg\Z
 \tag{3}
\]
tiles $\Z$ with every $A_i$.

\medskip
\noindent\textbf{The case $p=4$.}
Choose one integer spectrum $A_0/4$, translate $A_0$ so that $0\in A_0$,
and form the dephased Hadamard matrix with rows $v(a)$, $a\in A_0$.
The classification of $4\times4$ complex Hadamard matrices says that,
after row and column permutations, its rows have the form
\[
 \mathbf1,\quad u,\quad w,\quad uw,
 \qquad u^2=\mathbf1,\quad u\ne\mathbf1;
 \tag{4}
\]
this is the usual one-parameter family $F_4^{(1)}(t)$ recalled in
Herr--Wiegand~\cite{herrwiegand}.  Since every row lies in the cyclic
orbit $v(\Z)$, (4) gives a nonzero integer in the kernel of $v$.
Write
\[
 \ker v=q\Z.
\]
The image is cyclic of order $q$, and its unique element of order two is
$v(q/2)$.  Also
\[
 \ell_j:=\frac{q\lambda_j}{4}\in\Z,
 \qquad
 v(n)=\left(e^{2\pi i n\ell_j/q}\right)_{j=0}^3.
 \tag{5}
\]
The residues $L=\{\ell_0,\ldots,\ell_3\}\subset\Z/q\Z$ are distinct.

Apply (4) to each $A_i$ after translating one of its elements to zero.
Uniqueness of the order-two element in a cyclic group shows that, modulo
$q$, $A_i$ is a translate of
\[
 J(e_i)=\{0,N,e_i,e_i+N\},\qquad N=q/2.
 \tag{6}
\]
We next show that all $e_i$ have the same 2-adic valuation.  Orthogonality
of the rows $v(0)$ and $v(N)$ says that $L$ has two even and two odd
elements.  Fix the two even representatives $\ell',\ell''$ and put
$\delta=\ell'-\ell''$, an even nonzero integer modulo $q$.  For a row
set (6), orthogonality of $v(e_i)$ to both $v(0)$ and $v(N)$ gives
\[
 \sum_{\ell\in L}e^{2\pi i e_i\ell/q}=0,
 \qquad
 \sum_{\ell\in L}(-1)^\ell e^{2\pi i e_i\ell/q}=0.
\]
Adding the equations shows that the two even terms cancel, hence
\[
 e_i\delta\equiv q/2\pmod q.
 \tag{7}
\]
Because $\delta$ is even, (7) first implies $4\mid q$.  If
$s=\nu_2(q)$, then writing (7) as
$e_i\delta=(q/2)(2k+1)$ yields
\[
 \nu_2(e_i)=s-1-\nu_2(\delta)=:t<s-1,
 \tag{8}
\]
independently of $i$.

Set $d=\gcd(N,e_1,\ldots,e_m)$.  Equation (8) says that $N/d$ is even
and every $e_i/d$ is odd.  In $\Z/N\Z$ define
\[
 C=\{r+2dk:0\le r<d,\ 0\le k<N/(2d)\}.
 \tag{9}
\]
Translation by any $e_i$ interchanges $C$ and its complement, so
\(\{0,e_i\}\oplus C=\Z/N\Z\).  Viewing $C$ in $\Z/q\Z$, (6) gives
\[
 J(e_i)\oplus C=\Z/q\Z
\]
for every $i$.  Translates of $J(e_i)$ have the same complement.
Consequently $C+q\Z$ is a common tiling set for all the original $A_i$.

\medskip
\noindent\textbf{The case $p=5$.}
Again start with one translated spectrum and its dephased Hadamard
matrix.  Haagerup's order-five classification, in the explicit form of
Matolcsi--Ruzsa--Weiner~\cite[Proposition~2.1]{mrw}, says that every
$5\times5$ complex Hadamard matrix is equivalent to the Fourier matrix
$F_5$.  After dephasing, its rows are, up to permutations,
\[
 \mathbf1,r,r^2,r^3,r^4,qquad r^5=\mathbf1,quad r\ne\mathbf1.
 \tag{10}
\]
Because these rows lie in $v(\Z)$, (10) again makes $\ker v=q\Z$ for
some $q$.  The cyclic image has an element of order five, so $5\mid q$,
and it has a unique subgroup of order five.  Applying the same
classification to every $A_i$ shows that, modulo $q$, each $A_i$ is a
translate of
\[
 \{0,q/5,2q/5,3q/5,4q/5\}.
\]
Thus the single set
\[
 \{0,1,\ldots,q/5-1\}+q\Z
\]
is a common tiling complement.  This completes all five cases.
\end{proof}

\section{Consequence for spectral subsets of the line}

\begin{corollary}\label{line}
Let $\Omega\subset\R$ be bounded, measurable, spectral, and
$|\Omega|=1$.  If $\Omega$ has a spectrum
\[
 \Lambda=\{\lambda_0=0,\ldots,\lambda_{p-1}\}+p\Z
\]
with $p\le5$, then $\Omega$ tiles $\R$ by translations.  A tiling set
may be chosen inside $p^{-1}\Z$.
\end{corollary}

\begin{proof}
Proposition~2.8 of the source~\cite{source} decomposes $\Omega$ into
$p$-point integer fibers $\Omega_x$ which, for almost every $x$, have
the common spectrum
$p^{-1}\{\lambda_0,\ldots,\lambda_{p-1}\}$.  Boundedness leaves only
finitely many distinct fibers.  Theorem~\ref{main} gives them a common
integer tiling complement.  The fiber reconstruction in the same
proposition then tiles $\Omega$ by the corresponding subset of
$p^{-1}\Z$.
\end{proof}

\section{Verification, novelty, and limitations}

The finite-cyclic audit in \texttt{code/verify\_low\_period.py} enumerated
25,677 primitive column sets over the ranges
\[
 q\le60,36,24,18\quad\text{for}\quad p=2,3,4,5,
\]
respectively.  It found 1,527 nonempty spectrum families containing
1,677 normalized spectra and verified the common complements above in
every case.  In particular, the $p=4$ audit includes fixed column sets
with more than one row spectrum.  This computation is only a bounded
check; the proof is the argument above.

The local run indexes and bounded searches through 12 August 2026 used
the source title and arXiv id, ``Universal Tiling Conjecture'' with
$p=2,3,4,5$, low-order Hadamard submatrices, common spectra, and common
tiling complements.  No statement of $\mathrm{UTC}(p)$ for all
$p\le5$ was found.  Herr and Wiegand~\cite{herrwiegand} classify the
size-two and size-three Fourier-Hadamard compatibility data and end by
calling the relation between a row primitive set and its tiling set a
future direction; their paper supplies nearby context, not the common
complement theorem proved here.  Recent literature still describes the
one-dimensional general conjecture as open.

The method has a structural boundary at $p=6$: zero sums of six unit
vectors and order-six complex Hadamard matrices are non-rigid, and a
complete order-six classification is unavailable.  No claim is made for
$p\ge6$ or for the full one-dimensional Fuglede conjecture.  Human review
should focus on the dephasing-to-cyclic-subgroup step in the $p=4,5$
cases and on the common quotient complement (9).

\begin{thebibliography}{3}
\bibitem{source}
D.~E.~Dutkay and P.~E.~T.~Jorgensen,
\emph{On the universal tiling conjecture in dimension one},
J. Fourier Anal. Appl. 19 (2013), 467--477; arXiv:1209.0666,
\href{https://doi.org/10.1007/s00041-013-9264-7}{doi:10.1007/s00041-013-9264-7}.

\bibitem{herrwiegand}
J.~E.~Herr and T.~M.~Wiegand,
\emph{Identifying complex Hadamard submatrices of the Fourier matrices
via primitive sets}, Linear Algebra Appl. 617 (2021), 121--150;
arXiv:1909.13145,
\href{https://doi.org/10.1016/j.laa.2021.01.017}{doi:10.1016/j.laa.2021.01.017}.

\bibitem{mrw}
M.~Matolcsi, I.~Z.~Ruzsa, and M.~Weiner,
\emph{Real and complex unbiased Hadamard matrices}, arXiv:1201.0631,
2012.
\end{thebibliography}

\end{document}
